% Independent audit: path_stability passed the geometric proof and its refinements.
% Exact checks: check_high_path_cycle_layers.py; high-path-cycle-layers-results.json.
\section{Few internal cycles at fixed terminal-path density}
\label{sec:high-path-cycle-tradeoff}

A \emph{normalized two-pole} is a connected loopless multigraph with two distinct degree-two terminals.
Every other vertex has degree three.
Its rank $b$ and order $n$ satisfy $n=2b$.
Parallel edges remain distinct.

\begin{lemma}[Counting cycles through thin distance layers]
\label{lem:thin-layer-cycle-count}
Let $N$ be a bridgeless normalized two-pole of rank $b$, with terminals $s,t$.
Choose a shortest terminal path $A$, of length $\ell$.
Put
\[
 \Delta=b-\ell,\qquad H=\beta(N-V(A)).
\]
Then $\Delta\ge0$, and
\[
 c(N)\le 2^{19\Delta+12H+24+b/2}.
\]
If $N$ is simple, the term $12H$ can be omitted.
\end{lemma}

\begin{proof}
A connected bridgeless graph of maximum degree three has no cut vertex.
Let $L_i$ be its distance-$i$ layer from $s$.
For $1\le i<\ell$, the layer $L_i$ separates $s$ from $t$.
Thus $|L_i|\ge2$.
The endpoint layers each contain at least one vertex.
Consequently, $2b=|V(N)|\ge2\ell$, proving $\Delta\ge0$.

Mark all vertices outside the internal layers $L_1,\ldots,L_{\ell-1}$.
Also mark every internal layer containing at least three vertices.
The total excess above two vertices per internal layer, and one per endpoint layer, is $2\Delta$.
For $w\ge3$, one has $w\le3(w-2)$.
It follows that the marked set has at most $6\Delta+2$ vertices.
Every remaining layer has exactly two vertices, including its vertex of $A$.

Each vertex of an unmarked layer has a neighbor in the preceding layer and in the following layer.
Otherwise, deleting the other vertex of its layer would separate $s$ from $t$.
In particular, each such vertex has at most two incidences towards either neighboring layer.

If $N$ has parallel edges, mark every layer containing an endpoint of a parallel pair outside $V(A)$.
There are at most $H$ such pairs.
Indeed, their digons have disjoint edge supports and are independent in the cycle space of $N-V(A)$.
Parallel pairs have disjoint vertex sets, since the maximum degree is three.
A triple edge would form an isolated cubic component, which is impossible here.
No parallel pair has just one endpoint on $A$, because each vertex of $A$ has at most one incidence outside $A$.
Every parallel pair with both endpoints on $A$ joins consecutive path vertices, by geodesicity.
Each newly marked layer had two vertices, so at most $4H$ vertices are added.
The marked set now has at most $6\Delta+4H+2$ vertices.

Fix a cycle $C$ and its minimum and maximum distance layers $i\le j$.
Mark their vertices as well, unless their layers were already marked.
This adds at most four vertices.
Write $B$ for the final marked set, so
\[
 |B|\le6\Delta+4H+6.
\]
Fix all edges of $C$ incident with $B$.
There are at most $2^{3|B|}$ choices.

Call each unmarked layer \emph{ordinary}.
Consider an ordinary layer strictly between the two extreme layers of $C$.
The cycle must use both vertices of that layer.
Each vertex has one cycle edge to the preceding layer and one to the following layer.
To see this, every prefix cut meets a cycle evenly.
A layer with two vertices therefore carries two strands.
A same-layer edge would leave an odd prefix crossing.
Two separate caps at that layer would disconnect the cycle into a left part and a right part.
Ordinary layers outside the two extreme layers contribute no cycle edges.

Thus, across consecutive ordinary layers used by $C$, its edges form a perfect matching.
Such a cut has at most four edges.
For a simple cut, its perfect matching is unique unless the cut is $K_{2,2}$.
That cut has two perfect matchings.
The only other two-choice cut consists of a doubled edge of $A$ and one edge on the other strand.
Indeed, a doubled cross-strand edge would give an $A$ vertex two incidences outside $A$.
A doubled edge on the other strand has a marked endpoint.
A four-edge cut containing parallel edges must contain two doubled pairs, so it also has a marked endpoint.

Two cuts with two choices cannot be consecutive.
A $K_{2,2}$ cut already uses two incidences at each adjacent vertex.
An adjacent two-choice cut would force degree at least four at an $A$ vertex.
Two consecutive doubled edges of $A$ also force degree four.
Therefore each two-choice cut uses a disjoint pair of two-vertex layers.
There are at most $j-i+1$ such layers, and hence at most $(j-i+1)/2$ two-choice cuts.

The extreme layers, the edges incident with $B$, and these matching choices determine $C$ completely.
Let $D$ be the largest distance layer in $N$.
Every layer beyond $\ell$ consumes at least one of the $2\Delta$ excess vertices counted above.
Thus $D\le\ell+2\Delta\le b+2\Delta$.
The geometric interval sum satisfies
\[
 \sum_{0\le i\le j\le D}2^{(j-i+1)/2}
 \le \frac{\sqrt2}{(1-1/\sqrt2)^2}\,2^{D/2}
 <17\,2^{D/2}.
\]
Consequently,
\[
 c(N)<17\,2^{18\Delta+12H+18+(b+2\Delta)/2}
 \le2^{19\Delta+12H+24+b/2}.
\]
For a simple graph, no extra markings for parallel pairs are needed.
\end{proof}

\begin{theorem}[A fixed-density cycle tradeoff]
\label{thm:high-path-cycle-tradeoff}
For every $0<\varepsilon\le1/2$, there is a constant $K_\varepsilon$ such that every normalized two-pole $N$ of rank $b$ satisfies
\[
 P_{st}(N)\ge(1/2+\varepsilon)2^b
 \quad\Longrightarrow\quad
 c(N)\le K_\varepsilon 2^{b/2}.
\]
One may take $K_\varepsilon=2^{31C_\varepsilon+25}$, where $C_\varepsilon$ is the additive-stability constant.
For simple bridgeless two-poles, $2^{19C_\varepsilon+24}$ suffices.
\end{theorem}

\begin{proof}
First suppose that $N$ is bridgeless.
The additive path-stability proof gives
\[
 b-\ell\le C_\varepsilon,
 \qquad \beta(N-V(A))\le C_\varepsilon
\]
for a shortest terminal path $A$.
This application also permits a normalized multigraph.
Subdivide parallel edges to obtain a simple graph before applying the reduction.
Its suppression returns exactly $N$, and no cubic expansion is needed.
Thus the proof's residual-rank bound applies to this same $N$ and $A$.
Lemma~\ref{lem:thin-layer-cycle-count} proves the asserted estimate.

For a general normalized two-pole, every cyclic block lies on the terminal block route.
An off-route cyclic block would reduce the normalized path count to at most one half.
Every bridge therefore separates the terminals.
No bridge is incident with an original terminal, since its terminal shore would contain an off-route cycle.
Deleting all bridges leaves normalized bridgeless two-poles with distinct attachment terminals.
If two attachments coincided, the attached cyclic part would lie off the terminal route.
Every bridge shore contains a cycle, by the cubic degree count.

Write their ranks as $b_1,\ldots,b_q$.
Ranks add, terminal-path counts multiply, and cycle counts add.
Each factor's normalized path density is at least the product density, hence at least $1/2+\varepsilon$.
Apply the bridgeless bound, with constant $K_\varepsilon/2$, to every factor.
Since $b_i\ge1$ and $\sum_i b_i=b$, one has
\[
 \frac{\sum_i2^{b_i/2}}{2^{b/2}}
 \le q2^{-(q-1)/2}\le2.
\]
Summing the internal cycle counts proves the stated bound.
\end{proof}

\begin{corollary}[Simple cubic graphs with a two-edge cut]
\label{cor:simple-cubic-two-cut-error}
There is an absolute constant $K$ such that every connected bridgeless simple cubic graph $G$ of rank $r$ with a two-edge cut satisfies
\[
 c(G)\le2^{r-1}+K2^{r/2}.
\]
One may take $K=K_{1/4}$ from Theorem~\ref{thm:high-path-cycle-tradeoff}.
\end{corollary}

\begin{proof}
Deleting the cut leaves connected shores $N_1,N_2$.
The two attachment vertices in each shore are distinct.
Otherwise the shore's remaining incident edge would be a bridge of $G$.
Thus the shores are simple normalized two-poles.
Write $b_i=\beta(N_i)$, $p_i=P(N_i)$, and $c_i=c(N_i)$.
Then $b_i\ge2$, since a simple normalized two-pole cannot have rank one.
The exact decomposition is
\[
 r=b_1+b_2+1,\qquad c(G)=c_1+c_2+p_1p_2.
\]
Put $f_i=p_i/2^{b_i}$.
Call a shore high if $f_i\ge3/4$.

If both shores are high, the product term is at most $2^{r-1}$.
Theorem~\ref{thm:high-path-cycle-tradeoff} bounds the remaining sum by $K2^{r/2}$, since
$2^{b_1/2}+2^{b_2/2}\le2^{(b_1+b_2+1)/2}$.

Suppose only $N_2$ is high.
The crossing product has deficit at least $2^{b_1+b_2}/4$ from $2^{r-1}$.
Since $b_2\ge2$, this deficit is at least $2^{b_1}$.
It absorbs the trivial bound $c_1\le2^{b_1}-1$.
The remaining excess is at most $c_2$, which has the required bound.

Finally, suppose neither shore is high.
The crossing deficit is at least $7\cdot2^{b_1+b_2}/16$.
If $\max(b_1,b_2)\ge3$, then
\[
 2^{-b_1}+2^{-b_2}\le3/8<7/16.
\]
The two trivial internal cycle bounds are therefore absorbed, and $c(G)<2^{r-1}$.
The remaining case has $b_1=b_2=2$ and $r=5$.
The trivial cycle bound is already smaller than the asserted estimate.
\end{proof}

The final corollary addresses a broad decomposable subclass.
It does not prove the unrestricted half-exponent bound.
In particular, no two-edge cut is available in a cyclically three-edge-connected cubic kernel.

\begin{corollary}[All cubic two-edge cuts]
\label{cor:cubic-multigraph-two-cut-error}
There is an absolute constant $K'$ such that every connected loopless cubic multigraph with a two-edge cut satisfies
\[
 c(G)\le2^{r-1}+K'2^{r/2}.
\]
The same estimate holds for every connected cubic graph with a bridge.
\end{corollary}

\begin{proof}
First suppose that $G$ has a bridge.
Each shore contains a cycle, by the cubic degree count.
Their positive ranks $r_1,r_2$ add to $r$.
Thus
\[
 c(G)\le(2^{r_1}-1)+(2^{r_2}-1)\le2^{r-1}.
\]

Otherwise use the same two-pole decomposition as above.
Call a shore high when its normalized path density is at least $2/3$.
If both shores are high, apply Theorem~\ref{thm:high-path-cycle-tradeoff} with $\varepsilon=1/6$.
If neither is high, both ranks are at least two: a rank-one normalized two-pole is a saturated digon.
The crossing deficit is at least $(5/9)2^{b_1+b_2}$.
It absorbs both trivial internal bounds, because $2^{-b_1}+2^{-b_2}\le1/2$.

Suppose exactly one shore is high.
If its rank is at least two, its crossing deficit absorbs the low shore's internal cycles.
Indeed, $2^{b_{\rm high}}/3\ge1$.
The high shore's remaining cycles have the required half-exponent bound.

It remains to consider a high rank-one shore, which is a digon.
Let $N$ be the low shore, of rank $b$.
The exact formula gives
\[
 c(G)-2^{r-1}
 =1+c(N)+2P(N)-2^{b+1}
 \le1+c(N)-\frac23\,2^b.
\]
The asymptotic Erd\H{o}s--Simonovits bound supplies an absolute $b_0$ such that
$c(N)\le(2/3)2^b$ whenever $b\ge b_0$.
Thus the excess is at most one in this range.
For $b<b_0$, increase $K'$ to cover the trivial bound at the finitely many possible ranks.
\end{proof}

Except for the two-vertex triple-edge graph, a parallel pair in a bridgeless cubic multigraph gives a two-edge cut.
Consequently, the remaining cubic kernels for an unrestricted half-exponent bound are simple and three-edge-connected.
This reduction does not resolve that remaining class.

The exact checker reconstructs 878719 cycles on 11535 bridgeless two-poles.
Its input includes edge deletions from all 621 connected simple cubic graphs through fourteen vertices.
It also checks 94 capped-ladder and parallel-rail examples.
These checks verify the layer encoding and matching choices, rather than only the numerically loose final bound.
